% ============================================================
% Section 5.5: General Discussion and Interpretation of Results
% ============================================================
\section{General Discussion and Interpretation of Results}
\label{sec:general_discussion}

Based on the experimental results, the integrated system demonstrated high efficiency
in achieving its primary objectives: providing independent navigation and effective
virtual communication using electrooculography (EOG) signals. This performance stems
from key engineering and design choices, discussed and interpreted below.

% -----------------------------------------------------------
\subsection{Impact of Hybrid Signal Processing}
\label{subsec:impact_hybrid_signal_processing}

The results indicate that relying solely on digital filtering is insufficient for
processing microvolt-level EOG signals. Combining the AD620 instrumentation amplifier
with hardware analog filters was crucial for eliminating initial motion artifacts and
baseline instability. Subsequently, MATLAB digital filters (a $50\,\text{Hz}$ notch
filter and a moving average filter) suppressed power-line interference and smoothed
the waveform spikes. Distributing the processing workload between hardware and software
yielded a clean, stable signal while reducing the computational burden on the
microcontroller, enabling real-time command processing with minimal latency.

% -----------------------------------------------------------
\subsection{Overcoming Physiological Challenges}
\label{subsec:overcoming_physiological_challenges}

During initial testing, variations in skin impedance caused by perspiration, along with
involuntary facial muscle contractions (EMG artifacts), presented primary technical
challenges. These interferences manifested as baseline drift and biopotential waveform
distortions. To address these issues, a thresholding algorithm proved effective by
detecting voltage threshold crossings rather than relying on exact waveform matching.
This approach prevented false activations from rapid involuntary blinks, ensuring the
system responded exclusively to intentional, voluntary blinks and significantly reducing
false trigger rates.

% -----------------------------------------------------------
\subsection{Clinical Value of the Integrated Design}
\label{subsec:clinical_value_integrated_design}

Compared to existing assistive solutions for Locked-in Syndrome (LIS) and quadriplegic
patients, the proposed system offers improved user comfort and lower cognitive strain.
Unlike camera-based eye-tracking systems that require continuous visual focus---causing
severe eye fatigue and susceptibility to the ``Midas Touch'' problem---this design
provides asynchronous control. The user is not required to stare constantly at the
screen; instead, interaction occurs only through voluntary blinks when desired.
Furthermore, integrating the wheelchair and virtual keyboard into a unified software
platform reduces cognitive load by maintaining a consistent control modality across
both subsystems, with built-in standby states preventing unintended operational commands.

Overall, these findings demonstrate the clinical feasibility of the system as a
low-cost, effective assistive solution that provides functional autonomy while
maintaining user safety and operational reliability.
